\documentclass[10pt,a4paper]{amsart}
\usepackage{amsmath,amssymb,amsthm,graphicx}
% Retained preamble for the Stacks Project chapter morphisms.tex (Morphisms of Schemes).
% The chapter file pinned-original-source.tex carries only \input{preamble} on its first line, so no
% preamble text lives inside the chapter itself.  The Stacks Project uses one single shared preamble for
% every chapter; the snapshot stages it verbatim in sources/stacks-chow/preamble.tex.  The block below is
% copied mechanically, line for line, from that file's lines 14-18 (the Xy-pic setup) and lines 154-217
% (the native theorem declarations together with the mathematical macros).  Nothing was added, removed or
% redefined.  The wrapper class and the amsmath/amssymb/amsthm/graphicx base packages are supplied by the
% assembler, so the class/\externaldocument part of the shared preamble is not repeated here.
% For commutative diagrams we use Xy-pic
\usepackage[all]{xy}

% We use 2cell for 2-commutative diagrams.
\xyoption{2cell}
% Theorem environments.
%
\theoremstyle{plain}
\newtheorem{theorem}[subsection]{Theorem}
\newtheorem{proposition}[subsection]{Proposition}
\newtheorem{lemma}[subsection]{Lemma}

\theoremstyle{definition}
\newtheorem{definition}[subsection]{Definition}
\newtheorem{example}[subsection]{Example}
\newtheorem{exercise}[subsection]{Exercise}
\newtheorem{situation}[subsection]{Situation}

\theoremstyle{remark}
\newtheorem{remark}[subsection]{Remark}
\newtheorem{remarks}[subsection]{Remarks}

\numberwithin{equation}{subsection}

% Macros
%
\def\lim{\mathop{\mathrm{lim}}\nolimits}
\def\colim{\mathop{\mathrm{colim}}\nolimits}
\def\Spec{\mathop{\mathrm{Spec}}}
\def\Hom{\mathop{\mathrm{Hom}}\nolimits}
\def\Ext{\mathop{\mathrm{Ext}}\nolimits}
\def\SheafHom{\mathop{\mathcal{H}\!\mathit{om}}\nolimits}
\def\SheafExt{\mathop{\mathcal{E}\!\mathit{xt}}\nolimits}
\def\Sch{\mathit{Sch}}
\def\Mor{\mathop{\mathrm{Mor}}\nolimits}
\def\Ob{\mathop{\mathrm{Ob}}\nolimits}
\def\Sh{\mathop{\mathit{Sh}}\nolimits}
\def\NL{\mathop{N\!L}\nolimits}
\def\CH{\mathop{\mathrm{CH}}\nolimits}
\def\proetale{{pro\text{-}\acute{e}tale}}
\def\etale{{\acute{e}tale}}
\def\QCoh{\mathit{QCoh}}
\def\Ker{\mathop{\mathrm{Ker}}}
\def\Im{\mathop{\mathrm{Im}}}
\def\Coker{\mathop{\mathrm{Coker}}}
\def\Coim{\mathop{\mathrm{Coim}}}

% Boxtimes
%
\DeclareMathSymbol{\boxtimes}{\mathbin}{AMSa}{"02}

%
% Macros for moduli stacks/spaces
%
\def\QCohstack{\mathcal{QC}\!\mathit{oh}}
\def\Cohstack{\mathcal{C}\!\mathit{oh}}
\def\Spacesstack{\mathcal{S}\!\mathit{paces}}
\def\Quotfunctor{\mathrm{Quot}}
\def\Hilbfunctor{\mathrm{Hilb}}
\def\Curvesstack{\mathcal{C}\!\mathit{urves}}
\def\Polarizedstack{\mathcal{P}\!\mathit{olarized}}
\def\Complexesstack{\mathcal{C}\!\mathit{omplexes}}
% \Pic is the operator that assigns to X its picard group, usage \Pic(X)
% \Picardstack_{X/B} denotes the Picard stack of X over B
% \Picardfunctor_{X/B} denotes the Picard functor of X over B
\def\Pic{\mathop{\mathrm{Pic}}\nolimits}
\def\Picardstack{\mathcal{P}\!\mathit{ic}}
\def\Picardfunctor{\mathrm{Pic}}
\def\Deformationcategory{\mathcal{D}\!\mathit{ef}}

% --- appended 2026-10-12 (Chief): the project own definitions of the reference/slogan/history
% display environments, copied mechanically from the shared Stacks preamble lines 8-12
% (\usepackage{verbatim} plus the three \newenvironment lines, which the chapter preamble keeps
% verbatim from the project preamble). Nothing added, removed or redefined: these environments
% carry no mathematical content and the project itself typesets them as comments. Without these
% three lines a counted statement whose own text uses them fails with Environment slogan undefined.
\usepackage{verbatim}
\newenvironment{reference}{\comment}{\endcomment}
\newenvironment{slogan}{\comment}{\endcomment}
\newenvironment{history}{\comment}{\endcomment}
\usepackage{cleveref}
\crefname{originallemma}{original lemma}{original lemmas}
\Crefname{originallemma}{Original lemma}{Original lemmas}
\makeatletter
\newlabel{morphisms-lemma-characterize-relatively-ample}{{01VJ}{source}{}{sourcepointer-1}{}}
\newlabel{morphisms-lemma-characterize-relatively-ample@cref}{{[originallemma][0][]01VJ}{[1][0][]source}{}{}{}}
\makeatother
\crefname{originallemma}{original lemma}{original lemmas}
\Crefname{originallemma}{Original lemma}{Original lemmas}
\makeatletter
\newlabel{constructions-lemma-invertible-map-into-relative-proj}{{01O9}{source}{}{sourcepointer-2}{}}
\newlabel{constructions-lemma-invertible-map-into-relative-proj@cref}{{[originallemma][0][]01O9}{[1][0][]source}{}{}{}}
\makeatother
\crefname{originallemma}{original lemma}{original lemmas}
\Crefname{originallemma}{Original lemma}{Original lemmas}
\makeatletter
\newlabel{properties-lemma-ample-separated}{{09MP}{source}{}{sourcepointer-3}{}}
\newlabel{properties-lemma-ample-separated@cref}{{[originallemma][0][]09MP}{[1][0][]source}{}{}{}}
\makeatother
\crefname{originallemma}{original lemma}{original lemmas}
\Crefname{originallemma}{Original lemma}{Original lemmas}
\makeatletter
\newlabel{schemes-lemma-push-forward-quasi-coherent}{{01LC}{source}{}{sourcepointer-4}{}}
\newlabel{schemes-lemma-push-forward-quasi-coherent@cref}{{[originallemma][0][]01LC}{[1][0][]source}{}{}{}}
\makeatother
\crefname{originallemma}{original lemma}{original lemmas}
\Crefname{originallemma}{Original lemma}{Original lemmas}
\makeatletter
\newlabel{constructions-section-invertible-relative-proj}{{01O8}{source}{}{sourcepointer-5}{}}
\newlabel{constructions-section-invertible-relative-proj@cref}{{[originallemma][0][]01O8}{[1][0][]source}{}{}{}}
\makeatother
\crefname{originallemma}{original lemma}{original lemmas}
\Crefname{originallemma}{Original lemma}{Original lemmas}
\makeatletter
\newlabel{properties-lemma-map-into-proj}{{01PZ}{source}{}{sourcepointer-6}{}}
\newlabel{properties-lemma-map-into-proj@cref}{{[originallemma][0][]01PZ}{[1][0][]source}{}{}{}}
\makeatother
\crefname{originallemma}{original lemma}{original lemmas}
\Crefname{originallemma}{Original lemma}{Original lemmas}
\makeatletter
\newlabel{properties-lemma-ample-immersion-into-proj}{{01Q1}{source}{}{sourcepointer-7}{}}
\newlabel{properties-lemma-ample-immersion-into-proj@cref}{{[originallemma][0][]01Q1}{[1][0][]source}{}{}{}}
\makeatother
\crefname{originallemma}{original lemma}{original lemmas}
\Crefname{originallemma}{Original lemma}{Original lemmas}
\makeatletter
\newlabel{constructions-definition-relative-proj}{{01O0}{source}{}{sourcepointer-8}{}}
\newlabel{constructions-definition-relative-proj@cref}{{[originallemma][0][]01O0}{[1][0][]source}{}{}{}}
\makeatother
\crefname{originallemma}{original lemma}{original lemmas}
\Crefname{originallemma}{Original lemma}{Original lemmas}
\makeatletter
\newlabel{constructions-lemma-invertible-map-into-proj}{{01NK}{source}{}{sourcepointer-9}{}}
\newlabel{constructions-lemma-invertible-map-into-proj@cref}{{[originallemma][0][]01NK}{[1][0][]source}{}{}{}}
\makeatother
\crefname{originallemma}{original lemma}{original lemmas}
\Crefname{originallemma}{Original lemma}{Original lemmas}
\makeatletter
\newlabel{properties-lemma-open-in-proj-ample}{{01Q2}{source}{}{sourcepointer-10}{}}
\newlabel{properties-lemma-open-in-proj-ample@cref}{{[originallemma][0][]01Q2}{[1][0][]source}{}{}{}}
\makeatother
\crefname{originallemma}{original lemma}{original lemmas}
\Crefname{originallemma}{Original lemma}{Original lemmas}
\makeatletter
\newlabel{properties-lemma-ample-power-ample}{{01PT}{source}{}{sourcepointer-11}{}}
\newlabel{properties-lemma-ample-power-ample@cref}{{[originallemma][0][]01PT}{[1][0][]source}{}{}{}}
\makeatother
\crefname{originallemma}{original lemma}{original lemmas}
\Crefname{originallemma}{Original lemma}{Original lemmas}
\makeatletter
\newlabel{constructions-lemma-ample-on-proj}{{01MW}{source}{}{sourcepointer-12}{}}
\newlabel{constructions-lemma-ample-on-proj@cref}{{[originallemma][0][]01MW}{[1][0][]source}{}{}{}}
\makeatother
\crefname{originallemma}{original lemma}{original lemmas}
\Crefname{originallemma}{Original lemma}{Original lemmas}
\makeatletter
\newlabel{properties-lemma-ample-on-closed}{{01PU}{source}{}{sourcepointer-13}{}}
\newlabel{properties-lemma-ample-on-closed@cref}{{[originallemma][0][]01PU}{[1][0][]source}{}{}{}}
\makeatother
\crefname{originalreference}{original reference}{original references}
\Crefname{originalreference}{Original reference}{Original references}
\makeatletter
\newlabel{EGA}{{cite-EGA}{source}{}{sourcepointer-14}{}}
\newlabel{EGA@cref}{{[originalreference][0][]cite-EGA}{[1][0][]source}{}{}{}}
\makeatother
\title{Six equivalent formulations of relative ampleness}
\author{The Stacks Project Authors}
\date{stacks-project:morphisms.tex:local-snapshot-20261012; GNU FDL 1.2 or later}
\begin{document}
\maketitle
\setcounter{subsection}{1}
\noindent\textit{Editorial note.} All delivered mathematical text is the author's own, copied from the pinned Stacks Project chapter file morphisms.tex; no statement or proof step was rewritten, reordered or removed.\par\medskip
\noindent\textit{Editorial note.} The chapter file pinned-original-source.tex carries only \textbackslash{}input\{preamble\} on its first line, so no preamble text lives inside the chapter. The retained file sources/stacks-morphisms/retained-preamble.tex is a mechanical, line-for-line copy of the project's single shared preamble as staged in the snapshot. Nothing was added, removed or redefined, and the excerpt serves only as preamble, never as a span coordinate.\par\medskip
\noindent\textit{Editorial note.} Both counted spans start at the \textbackslash{}begin line of their environment, so the statement and the author's proof are delivered whole; surrounding prose of the section is deliberately outside the counted spans and is carried as standard prerequisites instead.\par\medskip
\noindent\textit{Editorial note.} The citation fields (source\_id, source\_version, authors, source\_url, license.name, source\_title) are identical to those of every other card this lane has produced from this chapter, so the batch selector treats them as one source.\par\medskip
\noindent\textit{Editorial note.} Screening note: the proof contains 'It is immediate from the definition that (1) implies (2) and (2) implies (3)', which is one legitimate implication inside a long argument; the author proves the remaining implications in full, so this is not a hint-only proof.\par\medskip
\noindent\textit{Allowed standard prerequisites (editorial).}\begin{itemize}
\item Ample invertible sheaves on a scheme over an affine base are separated, and the relative Proj construction. Properties, Lemma \textbackslash{}ref\{properties-lemma-ample-separated\}, Constructions, Section \textbackslash{}ref\{constructions-section-invertible-relative-proj\}, and Constructions, Lemma \textbackslash{}ref\{constructions-lemma-invertible-map-into-relative-proj\}, cited at proof lines 8719, 8735-8738, and 8743-8746. Provide the constructions used throughout.
\item Ampleness descends and ascends along the constructions used, and ample sheaves restrict to closed subschemes. Properties, Lemma \textbackslash{}ref\{properties-lemma-open-in-proj-ample\}, Properties, Lemma \textbackslash{}ref\{properties-lemma-ample-on-closed\}, and Properties, Lemma \textbackslash{}ref\{properties-lemma-ample-power-ample\}, cited at proof lines 8788-8795. Close the cycle at the end of the proof.
\end{itemize}\medskip
\ifdefined\hypertarget\hypertarget{sourcepointer-1}{}\fi
\noindent\textit{Source reference.} Original lemma 01VJ: focused-seeds/AGP/sources/stacks-snapshot-20261011/tags.txt, source lines 2448--2448. The lemma's own canonical tag, quoted so that readers can resolve the label inside the delivered text; it points at the counted statement itself and carries no proof content.\par\medskip
\ifdefined\hypertarget\hypertarget{sourcepointer-2}{}\fi
\noindent\textit{Source reference.} Original lemma 01O9: focused-seeds/AGP/sources/stacks-snapshot-20261011/tags.txt, source lines 2186--2186. Navigation pointer to a statement the counted proof refers to but does not contain; nothing load-bearing is delegated to it.\par\medskip
\ifdefined\hypertarget\hypertarget{sourcepointer-3}{}\fi
\noindent\textit{Source reference.} Original lemma 09MP: focused-seeds/AGP/sources/stacks-snapshot-20261011/tags.txt, source lines 12187--12187. Navigation pointer to a statement the counted proof refers to but does not contain; nothing load-bearing is delegated to it.\par\medskip
\ifdefined\hypertarget\hypertarget{sourcepointer-4}{}\fi
\noindent\textit{Source reference.} Original lemma 01LC: focused-seeds/AGP/sources/stacks-snapshot-20261011/tags.txt, source lines 2081--2081. Navigation pointer to a statement the counted proof refers to but does not contain; nothing load-bearing is delegated to it.\par\medskip
\ifdefined\hypertarget\hypertarget{sourcepointer-5}{}\fi
\noindent\textit{Source reference.} Original lemma 01O8: focused-seeds/AGP/sources/stacks-snapshot-20261011/tags.txt, source lines 2185--2185. Navigation pointer to a statement the counted proof refers to but does not contain; nothing load-bearing is delegated to it.\par\medskip
\ifdefined\hypertarget\hypertarget{sourcepointer-6}{}\fi
\noindent\textit{Source reference.} Original lemma 01PZ: focused-seeds/AGP/sources/stacks-snapshot-20261011/tags.txt, source lines 2248--2248. Navigation pointer to a statement the counted proof refers to but does not contain; nothing load-bearing is delegated to it.\par\medskip
\ifdefined\hypertarget\hypertarget{sourcepointer-7}{}\fi
\noindent\textit{Source reference.} Original lemma 01Q1: focused-seeds/AGP/sources/stacks-snapshot-20261011/tags.txt, source lines 2250--2250. Navigation pointer to a statement the counted proof refers to but does not contain; nothing load-bearing is delegated to it.\par\medskip
\ifdefined\hypertarget\hypertarget{sourcepointer-8}{}\fi
\noindent\textit{Source reference.} Original lemma 01O0: focused-seeds/AGP/sources/stacks-snapshot-20261011/tags.txt, source lines 2177--2177. Navigation pointer to a statement the counted proof refers to but does not contain; nothing load-bearing is delegated to it.\par\medskip
\ifdefined\hypertarget\hypertarget{sourcepointer-9}{}\fi
\noindent\textit{Source reference.} Original lemma 01NK: focused-seeds/AGP/sources/stacks-snapshot-20261011/tags.txt, source lines 2161--2161. Navigation pointer to a statement the counted proof refers to but does not contain; nothing load-bearing is delegated to it.\par\medskip
\ifdefined\hypertarget\hypertarget{sourcepointer-10}{}\fi
\noindent\textit{Source reference.} Original lemma 01Q2: focused-seeds/AGP/sources/stacks-snapshot-20261011/tags.txt, source lines 2251--2251. Navigation pointer to a statement the counted proof refers to but does not contain; nothing load-bearing is delegated to it.\par\medskip
\ifdefined\hypertarget\hypertarget{sourcepointer-11}{}\fi
\noindent\textit{Source reference.} Original lemma 01PT: focused-seeds/AGP/sources/stacks-snapshot-20261011/tags.txt, source lines 2242--2242. Navigation pointer to a statement the counted proof refers to but does not contain; nothing load-bearing is delegated to it.\par\medskip
\ifdefined\hypertarget\hypertarget{sourcepointer-12}{}\fi
\noindent\textit{Source reference.} Original lemma 01MW: focused-seeds/AGP/sources/stacks-snapshot-20261011/tags.txt, source lines 2137--2137. Navigation pointer to a statement the counted proof refers to but does not contain; nothing load-bearing is delegated to it.\par\medskip
\ifdefined\hypertarget\hypertarget{sourcepointer-13}{}\fi
\noindent\textit{Source reference.} Original lemma 01PU: focused-seeds/AGP/sources/stacks-snapshot-20261011/tags.txt, source lines 2243--2243. Navigation pointer to a statement the counted proof refers to but does not contain; nothing load-bearing is delegated to it.\par\medskip
\ifdefined\hypertarget\hypertarget{sourcepointer-14}{}\fi
\noindent\textit{Source reference.} Original reference cite-EGA: sources/stacks-morphisms/pinned-original-source.tex, source lines 8679--8679. Bibliographic pointer appearing in the counted text; it carries no mathematical argument.\par\medskip
\begin{lemma}
\label{lemma-characterize-relatively-ample}
\begin{reference}
\cite[II, Proposition 4.6.3]{EGA}
\end{reference}
Let $f : X \to S$ be a quasi-compact morphism of schemes.
Let $\mathcal{L}$ be an invertible sheaf on $X$.
The following are equivalent:
\begin{enumerate}
\item The invertible sheaf $\mathcal{L}$ is $f$-ample.
\item There exists an open covering $S = \bigcup V_i$
such that each $\mathcal{L}|_{f^{-1}(V_i)}$ is ample
relative to $f^{-1}(V_i) \to V_i$.
\item There exists an affine open covering $S = \bigcup V_i$
such that each $\mathcal{L}|_{f^{-1}(V_i)}$ is ample.
\item There exists a quasi-coherent graded $\mathcal{O}_S$-algebra
$\mathcal{A}$ and a map of graded $\mathcal{O}_X$-algebras
$\psi : f^*\mathcal{A} \to \bigoplus_{d \geq 0} \mathcal{L}^{\otimes d}$
such that $U(\psi) = X$ and
$$
r_{\mathcal{L}, \psi} :
X
\longrightarrow
\underline{\text{Proj}}_S(\mathcal{A})
$$
is an open immersion (see Constructions, Lemma
\ref{constructions-lemma-invertible-map-into-relative-proj} for notation).
\item The morphism $f$ is quasi-separated and
part (4) above holds with
$\mathcal{A} = f_*(\bigoplus_{d \geq 0} \mathcal{L}^{\otimes d})$
and $\psi$ the adjunction mapping.
\item Same as (4) but just requiring $r_{\mathcal{L}, \psi}$
to be an immersion.
\end{enumerate}
\end{lemma}

\begin{proof}
It is immediate from the definition that (1) implies (2) and
(2) implies (3). It is clear that (5) implies (4).

\medskip\noindent
Assume (3) holds for the affine open covering $S = \bigcup V_i$.
We are going to show (5) holds.
Since each $f^{-1}(V_i)$ has an ample invertible sheaf we see
that $f^{-1}(V_i)$ is separated
(Properties, Lemma \ref{properties-lemma-ample-separated}).
Hence $f$ is separated. By
Schemes, Lemma \ref{schemes-lemma-push-forward-quasi-coherent}
we see that $\mathcal{A} = f_*(\bigoplus_{d \geq 0} \mathcal{L}^{\otimes d})$
is a quasi-coherent graded $\mathcal{O}_S$-algebra.
Denote $\psi : f^*\mathcal{A} \to \bigoplus_{d \geq 0} \mathcal{L}^{\otimes d}$
the adjunction mapping.
The description of the open $U(\psi)$ in
Constructions, Section
\ref{constructions-section-invertible-relative-proj}
and the definition of ampleness
of $\mathcal{L}|_{f^{-1}(V_i)}$ show that $U(\psi) = X$.
Moreover, Constructions,
Lemma \ref{constructions-lemma-invertible-map-into-relative-proj} part (3)
shows that the restriction of $r_{\mathcal{L}, \psi}$ to
$f^{-1}(V_i)$ is the same as the morphism from
Properties, Lemma \ref{properties-lemma-map-into-proj}
which is an open immersion according to
Properties, Lemma \ref{properties-lemma-ample-immersion-into-proj}.
Hence (5) holds.

\medskip\noindent
Let us show that (4) implies (1). Assume (4).
Denote $\pi : \underline{\text{Proj}}_S(\mathcal{A}) \to S$
the structure morphism. Choose $V \subset S$ affine open. By
Constructions, Definition \ref{constructions-definition-relative-proj}
we see that $\pi^{-1}(V) \subset \underline{\text{Proj}}_S(\mathcal{A})$
is equal to $\text{Proj}(A)$ where $A = \mathcal{A}(V)$
as a graded ring. Hence $r_{\mathcal{L}, \psi}$ maps
$f^{-1}(V)$ isomorphically onto
a quasi-compact open of $\text{Proj}(A)$.
Moreover, $\mathcal{L}^{\otimes d}$ is isomorphic to the pullback of
$\mathcal{O}_{\text{Proj}(A)}(d)$ for some $d \geq 1$.
(See part (3) of Constructions,
Lemma \ref{constructions-lemma-invertible-map-into-relative-proj}
and the final statement of Constructions,
Lemma \ref{constructions-lemma-invertible-map-into-proj}.)
This implies that $\mathcal{L}|_{f^{-1}(V)}$ is ample by
Properties, Lemmas \ref{properties-lemma-open-in-proj-ample}
and \ref{properties-lemma-ample-power-ample}.

\medskip\noindent
Assume (6). By the equivalence of (1) - (5) above we see that the
property of being relatively ample on $X/S$ is local on $S$. Hence
we may assume that $S$ is affine, and we have to show that
$\mathcal{L}$ is ample on $X$. In this case the morphism
$r_{\mathcal{L}, \psi}$ is identified with the morphism, also denoted
$r_{\mathcal{L}, \psi} : X \to \text{Proj}(A)$ associated to the map
$\psi : A = \mathcal{A}(V) \to \Gamma_*(X, \mathcal{L})$.
(See references above.) As above we also see that
$\mathcal{L}^{\otimes d}$ is the pullback of the sheaf
$\mathcal{O}_{\text{Proj}(A)}(d)$ for some $d \geq 1$.
Moreover, since $X$ is quasi-compact we
see that $X$ gets identified with a closed subscheme of a
quasi-compact open subscheme $Y \subset \text{Proj}(A)$.
By
Constructions, Lemma
\ref{constructions-lemma-ample-on-proj}
(see also
Properties, Lemma
\ref{properties-lemma-open-in-proj-ample})
we see that $\mathcal{O}_Y(d')$ is an ample invertible sheaf on
$Y$ for some $d' \geq 1$. Since the restriction of an ample
sheaf to a closed subscheme is ample, see
Properties, Lemma
\ref{properties-lemma-ample-on-closed}
we conclude that the pullback of
$\mathcal{O}_Y(d')$ is ample. Combining these results with
Properties, Lemma
\ref{properties-lemma-ample-power-ample}
we conclude that $\mathcal{L}$ is ample as desired.
\end{proof}

\end{document}
